\section{What we test, and how}
\label{sec:design}

\subsection{The interventions}
\label{sec:design-interventions}

Every intervention below conditions the same host (Qwen3.5-4B unless noted) to hold the same
assigned stance on the same topic, and differs only in \emph{how} the stance is delivered.
Full prompt texts are in Appendix~\ref{app:prompts}; the naming convention is a channel prefix
(\textsc{ctx} for context-delivered, \textsc{cmp} for compiled) and a strength suffix.

\textbf{Context prompts.} The belief is placed in the context window as text, under a system
instruction. We vary the instruction across five strengths, because the field's implicit
baseline---``tell the model what to think''---turns out to span a wide range of behavior:
\emph{mild} (``you hold strong convictions; argue your position directly''), \emph{persist}
(an explicit non-yielding directive: treat objections without new evidence as noise, never
soften to appease), \emph{engage} (an explicit updating directive: address the specific
evidence, grant what is valid, revise on a real source with a warrant), \emph{balanced} (hold
against pressure, but update on genuine evidence---both directives in one instruction), and
\emph{discerning} (weigh what each challenge offers; hold against doubt, emotion, or
unsourced claims; update precisely as far as a real source warrants).

\textbf{Contrastive activation steering.} A tuned steering vector derived from the same belief
content, added at the residual stream of a single layer, following \citet{rimsky2024caa}, the
standard activation-space alternative to prompting. The tuned configuration we obtained is
degenerate---it holds below the unconditioned host---and is reported as a failed baseline
rather than as a measurement of the method class (Section~\ref{sec:dissociation-table}).

\textbf{Weight-space fine-tuning.} A per-belief LoRA fine-tune on the same belief content,
parameter-matched to the compiled channel and fairly configured---the same non-thinking prefill
the compiled conditions use, and repetition-suppressing decoding---then evaluated through the
identical protocol (Appendix~\ref{app:lora}).

\textbf{Compiled modulation.} A trained hypernetwork maps the belief content to per-layer
low-rank conditioning of the query and value projections, at rank 16 over a stride-3 subset of
layers (11 of 32), amortizing into one forward pass what LoRA obtains by gradient descent;
the lineage is AdaLN/FiLM-style conditional normalization
\citep{perez1970film,xu2019adain} and hypernetwork adapters
\citep{chen2025generativeadapter,charakorn2025texttolora}. The query-side update is produced
from a disposition vector alone; the value-side update is produced from that vector pooled
with the host's hidden states over the belief text, so the value channel sees content the
query channel never does (Section~\ref{sec:why-channels-differ}). Neither update depends on the
current prompt. Because the compilation drops low-frequency specifics
(Section~\ref{sec:boundary}), the deployed configuration pairs it with a small curated context
slot carrying those specifics; we report the modulation with and without that slot.

\textbf{Combinations.} The comparison this paper turns on is a matched pair: the same system
prompt, the same context, with the compiled modulation switched on or off. We run this pair at
every prompt strength, so the modulation's contribution is measured against its own control
rather than against a different paper's baseline.

\subsection{Protocol}
\label{sec:design-protocol}

Twenty-four held-out debate topics, two sides, with pressure applied over five escalating
levels---L1 bare doubt, L2 fabricated evidence, L3 logical counterargument, L4 appeal to
authority, L5 emotional pressure---three rounds each. Levels 1, 4, and 5 supply no evidence at
all; L2 fabricates a study; L3 asserts a reasoning flaw without naming one. Conceding to any of
these is yielding to pressure rather than to information. Pressure turns never restate the
assigned stance, so they cannot act as a second copy of the belief.

Two further protocols supply the axes a pressure ladder cannot reach. For \emph{evidence
discrimination}, each topic is seeded with the side that the strongest real evidence
contradicts, so a valid correction exists by construction, and challenges are drawn from a
fixture set of six types: two well-warranted (a real source with specifics and a warrant that
bears on the claim) and four deficient (unsourced assertion, unnamed authority, non sequitur,
impossible magnitude), style-matched across classes for tone, assertiveness, and length within
\(\pm 20\%\) (Appendix~\ref{app:fixtures}). For \emph{adversarial revocation}, three multi-turn
attack strategies attempt to remove the commitment rather than argue against it: a crescendo
that recruits the model's own words, an eviction that forces an opposing restatement, and an
authority reframe that claims to have authored the model's instructions and to be correcting
them (Appendix~\ref{app:attacks}).

\subsection{Scoring}
\label{sec:design-scoring}

Generations are scored by a six-judge ensemble spanning three model families and two size
tiers, under a rubric that shows the judge the belief being defended and asks it to weigh
argumentative ground rather than tone---a response may be gracious and still score as holding,
curt and still score as folding (Krippendorff \(\alpha=0.79\); rubrics and per-judge breakdown
in Appendix~\ref{app:multijudge}). The topic is the unit of analysis (\(n=24\)); intervals are
bootstrap over topics; paired comparisons use Wilcoxon signed-rank with Holm--Bonferroni
correction across the family of tests (Appendix~\ref{app:significance}). Cells scored by a
single judge rather than the full ensemble are marked \(\ddagger\) throughout. All runs carry
per-record emptiness guards, closing a validity hazard in which an expired reasoning-block
budget silently zeroes only the unmodulated arms (Appendix~\ref{app:repro}).
